% ==============================================================================
% PAGE 8: ABSTRACT
% ==============================================================================
\thispagestyle{empty}
\begin{tikzpicture}[remember picture, overlay]
  \draw[line width=2.0pt, color=Goldenrod!90!black]
    ([xshift=12mm, yshift=-12mm]current page.north west) rectangle ([xshift=-12mm, yshift=12mm]current page.south east);
  \draw[line width=0.8pt, color=Goldenrod!80!black]
    ([xshift=14.5mm, yshift=-14.5mm]current page.north west) rectangle ([xshift=-14.5mm, yshift=14.5mm]current page.south east);
\end{tikzpicture}

\vspace*{-5mm}
\begin{center}
  {\fontsize{18}{22}\selectfont \bfseries \textcolor{red!80!black}{\underline{ABSTRACT}}}\\[6mm]
\end{center}

\noindent
\begin{spacing}{1.25}
{\fontsize{10.2}{14}\selectfont
The increasing global energy demand and depletion of fossil fuels have accelerated the need for sustainable energy solutions. Solar energy, being clean, renewable, and abundant, is a reliable alternative for thermal and electrical applications. In this work, a solar water heater and steam generator system is designed and fabricated using a cylindrical parabolic concentrator (CPC). The CPC focuses incident solar radiation onto a receiver tube, thereby enhancing the thermal efficiency. Unlike flat plate collectors, the concentrator achieves higher temperatures suitable for both water heating and steam generation. The proposed design ensures optimum utilization of solar energy for domestic, industrial, and agricultural applications.\\[2.5mm]
The system incorporates real-time automated solar tracking to overcome limitations of fixed-angle collectors. A dual-axis tracking mechanism powered by sensors and microcontroller logic adjusts the concentrator to maintain perpendicularity to solar rays. This ensures maximum concentration of energy throughout the day, significantly improving thermal output. The tracking system reduces losses caused by diurnal variations and seasonal shifts in solar position. Automation also minimizes human intervention and enhances operational reliability. With improved efficiency, the system provides consistent steam production even during partial sunlight conditions.\\[2.5mm]
The fabricated prototype consists of a cylindrical concentrator, absorber tube, insulated water storage, and control electronics. Mild steel and reflective aluminum are used for the concentrator frame and surface to ensure durability and high reflectivity. The receiver tube, coated with selective absorptive material, minimizes heat losses while maximizing absorption. Heat transfer to the working fluid (water) occurs rapidly, raising its temperature to produce hot water and steam. Insulated storage prevents thermal losses and allows continuous supply even after sunset. The compact and modular design enables easy scalability for domestic as well as industrial usage.\\[2.5mm]
Experimental testing demonstrated the system's ability to achieve water heating up to boiling point and generate steam at controlled rates. Real-time tracking improved efficiency by 30--40\% compared to fixed concentrators. The study concludes that the developed system is an efficient, low-cost, and eco-friendly alternative to conventional water heating and steam generation methods. It offers significant potential in rural areas where grid energy is limited. Additionally, the system reduces carbon emissions, promoting green technology adoption.\\[2.5mm]
Further research may involve hybridization with other renewable systems and performance optimization under varying climatic conditions.
}
\end{spacing}
\clearpage
